\chapter*{常见缩写和术语先导}
\addcontentsline{toc}{chapter}{常见缩写和术语先导}

代码实践卷会把第二册的概念落到工程对象。第一次看到下面这些词时，先把它们理解成“代码里必须有边界的对象”，而不是抽象口号；每个对象都要能在代码、测试或报告里找到证据。

\begin{description}
  \item[reference] 参考实现。它可以慢，但必须清楚、确定、可复查，用来定义正确结果。
  \item[oracle] 判定器。它比较 reference 和待测实现，判断输出是否满足同一语义和误差合同。
  \item[contract] 合同。这里不是法律文件，而是代码边界的明确约定：输入是什么，输出是什么，错误怎样表达，哪些行为不承诺。
  \item[RecordId] 记录身份。它把一条输入记录定位到 input name、generation、byte offset 和 line number，方便诊断和重放。
  \item[diagnostic] 诊断记录。它不是普通日志，而是给测试、报告和修复定位使用的结构化错误事实。
  \item[checksum] 校验和。它把输出内容压成可比较摘要，用来检查结果是否稳定、文件是否被截断或误读。
  \item[CPU] Central Processing Unit，中央处理器。本卷的实验默认先落在本地 CPU 上，再逐步讨论 SIMD、线程和 I/O。
  \item[UI] User Interface，用户界面。第一章提到 UI 时，重点是错误和诊断如何被用户看懂，而不是先做图形界面。
  \item[C++20] C++ 标准版本名。本卷代码以 C++20 为默认边界，原因是它同时提供现代类型、span、同步原语和较好的工程表达能力。
  \item[CMake] C++ 工程常用构建系统。本卷用它组织 target、测试、benchmark 和工具程序。
  \item[\code{std::}] C++ 标准库命名空间前缀。本卷第一章里的 \code{std::vector}、\code{std::string}、\code{std::map}、\code{std::uint64_t} 都来自标准库；读代码时先问它们表达“拥有一组元素”“保存字节串”“按键查值”还是“固定宽度整数”。
  \item[\code{InputCase/ReferenceResult/ParseDiagnostic}] 本卷第一章的项目内示例名。InputCase 表示一条测试输入；ReferenceResult 表示参考实现的输出；ParseDiagnostic 表示解析失败时的结构化诊断。它们不是行业缩写，而是代码合同的名字。
  \item[SIMD] Single Instruction, Multiple Data，一条指令处理多个数据。实践卷后续会用它优化热循环。
  \item[I/O/EOF] I/O 是 Input/Output，输入输出；EOF 是 End Of File，表示输入结束。实践卷会把读取、短写、完成事件和业务提交分开处理。
  \item[CLI/CI] CLI 是命令行接口，便于脚本化运行实验；CI 是持续集成，用自动构建和测试保护项目边界。
  \item[JSON/CSV] JSON 适合保存结构化配置和 manifest；CSV 适合保存表格化 benchmark、trace 摘要和对比结果。
  \item[manifest] 机器可读清单，记录输入版本、输出位置、校验和、参数和提交事实。
  \item[checkpoint] 检查点，保存可恢复状态，让失败后可以从已提交边界继续。
  \item[attempt] 一次物理尝试。同一个逻辑任务可能因为失败重试而产生多个 attempt。
  \item[epoch] 版本或轮次编号，用来区分不同时刻的状态和提交边界。
  \item[MessageBus] 消息总线。本卷里指组件之间传递任务、完成事件和控制消息的工程边界。
  \item[batch/split] split 是输入切分后的逻辑片段；batch 是热路径里成批处理的一组记录。split 偏调度身份，batch 偏执行效率。
  \item[idempotent] 幂等。重复执行同一逻辑请求不会产生额外业务结果，是重试、manifest 和恢复的核心要求。
  \item[benchmark/trace] benchmark 是可重复的性能测量；trace 是按时间记录的事件流。实践卷里二者都要带输入版本、构建选项和结果校验。
  \item[target/span/event/generation] target 是 CMake 构建目标或业务目标；span 是一段连续数据的非拥有视图；event 是状态变化或完成事实；generation 是同一逻辑输入的新版本编号。
  \item[offset/line number] offset 是输入中的字节偏移，line number 是行号。诊断同时保存二者，是为了既能机器重放，也能让人快速定位原始文本。
  \item[工程对象名] \code{ManifestEntry} 表示清单中的一条输入或输出记录；\code{InputSplit} 表示输入被切出来的逻辑片段；\code{HotBatch} 表示热路径中被成批处理的数据；\code{QueueContract} 表示队列必须满足的容量、顺序、关闭和失败语义；\code{RunReport} 表示一次运行的结构化结果。
  \item[诊断字段名] \code{input\_name} 是输入来源名；\code{input\_generation} 是输入版本；\code{byte\_offset} 是字节位置；\code{line\_number} 是行号；\code{error\_kind} 是错误分类；\code{records\_total} 是处理记录总数。这些字段的目标是让失败能被复现，而不是让日志看起来复杂。
  \item[UTF-8] Unicode 的常见变长字节编码。实践卷保存 offset 时默认按字节定位；如果 UI 要按“第几个字符”显示，就必须额外处理 UTF-8 解码，不能把字节下标直接当字符下标。
\end{description}
